\documentclass[11pt]{article}

% =====================================================================
%  THEORY NOTES (backstage) for A-DCVAE  -- Revision 3
%  Self-contained mathematical note. NOT the paper. Kept fully rigorous;
%  the paper will abbreviate/summarize as needed. Reused as the formal
%  foundation for subsequent theorems.
%  Compile standalone:  pdflatex theory_notes_rev3.tex
%
%  Rev.2 incorporated a rigor review: (i) measurability of the kernel is
%  now proved; (ii) the coupling is formalized via an explicit joint
%  kernel + Ionescu--Tulcea; (iii) the deterministic-collapse remark is
%  corrected (Doeblin is sufficient, not necessary); (iv) the standard
%  facts used (TV-completeness, Leb(dK)=0, kernel TV-nonexpansiveness,
%  projection properties) are stated as cited/proved lemmas; plus assorted
%  fixes (epsilon<=1 via A=Omega, typos).
%
%  Rev.3 (this revision). Constraint: everything up to and including
%  Theorem A (Secs. 1--9) is left VERBATIM; only the preamble (cosmetic)
%  and Secs. 10--13 + bibliography are touched. Changes, each marked
%  inline with a [Rev.3-n] comment:
%   [Rev.3-1]  Cosmetic only: theorem/definition/remark styles now break
%              the line after the heading, and proofs start their body on
%              a fresh line. No mathematical content is affected.
%   [Rev.3-2]  New Fact (kernmeas): on a Polish target, the evaluation
%              sigma-algebra on P(Om) equals the Borel sigma-algebra of
%              the weak topology. This is the missing bridge that (a)
%              legitimizes applying Villani Cor. 5.22 to our kernels and
%              (b) makes mixtures of coupling kernels well defined.
%   [Rev.3-3]  Fact (msel) restated as (i)-(iii) with a full proof;
%              "equivalently" weakened to a proved "in particular"; the
%              operation \int gamma_s rho(ds) is now called MIXING (the
%              gluing lemma is a different operation) here and wherever
%              it is invoked (Lemma sweep, Lemma perturb).
%   [Rev.3-4]  Sec. 13 target paragraph: removed the category error
%              "Pi-dagger is K-dagger-invariant" (K-dagger acts on
%              P(Om_Theta), not on joint laws); the needed statement
%              nu-dagger K-dagger = nu-dagger is now PROVED by the
%              sweep-reproduction argument (same as Prop. inc), and the
%              kernel-ness of K-dagger is pointed to Fact compose.
%   [Rev.3-5]  Assumption citations of Lemma opdisc and Theorem B now
%              include Assumption suptrue (it was silently used: domain
%              of the true sweep, definition of eps_H/eps_P, and the
%              Lipschitz domain coverage all depend on it). suptrue is
%              kept in Sec. 13 (not moved to Sec. 6) so that Secs. 1--9
%              stay untouched; logically it is standing for Sec. 13.
%   [Rev.3-6]  Prop. inc: standing assumptions made explicit; Remark
%              pileftpi: kappa<1 noted as in force.
%   [Rev.3-7]  Remark edist: the law of Y (stationary pre-projection
%              proposal) is specified; sqrt(d) dimension factors flagged.
%   [Rev.3-8]  Remark nodouble: hedge sharpened -- the swap-learned-first
%              triangle route suggests the constant closes with
%              kappa, L_H, L_P alone (no regularity of the true
%              conditionals); precise constant still deferred.
%   [Rev.3-9]  Lemma condlip / Lemma sweep statements: quantifier domains
%              Om_Theta, Om_H made explicit (matching the domains in
%              Assumption lip).
%   [Rev.3-10] Notation in Secs. 10--13 unified to the bold convention of
%              Secs. 1--8 (\bth, \bm u, \bm z_H, ...). Purely notational.
%   [Rev.3-11] Bibliography parentheticals extended (Villani Cor. 5.22;
%              Kallenberg: measurable structure of P(Om)).
% =====================================================================

\usepackage[margin=1in]{geometry}
\usepackage{amsmath,amssymb,amsthm,mathtools}
\usepackage{bm}
\usepackage{hyperref}

% ---------------------------------------------------------------------
% [Rev.3-1] COSMETIC: break-after-heading theorem styles.
% WHY: run-in headings ("Theorem 4.1 (name). text...") reduce readability.
% HOW: custom amsthm styles with \newline as the head/body separator, so
% every statement body starts on its own line; the proof environment is
% redefined so "Proof." also stands on its own line. Display only ---
% numbering, counters, and all statement bodies are unchanged.
% ---------------------------------------------------------------------
\newtheoremstyle{plainbreak}%
  {10pt}{10pt}{\itshape}{}{\bfseries}{.}{\newline}{}
\newtheoremstyle{definitionbreak}%
  {10pt}{10pt}{\normalfont}{}{\bfseries}{.}{\newline}{}
\newtheoremstyle{remarkbreak}%
  {8pt}{8pt}{\normalfont}{}{\itshape}{.}{\newline}{}

\theoremstyle{definitionbreak}
\newtheorem{definition}{Definition}[section]
\newtheorem{assumption}{Assumption}
\newtheorem{notation}{Notation}[section]

\theoremstyle{plainbreak}
\newtheorem{theorem}{Theorem}[section]
\newtheorem{lemma}[theorem]{Lemma}
\newtheorem{proposition}[theorem]{Proposition}
\newtheorem{corollary}[theorem]{Corollary}
\newtheorem{fact}[theorem]{Fact}

\theoremstyle{remarkbreak}
\newtheorem{remark}{Remark}[section]

% [Rev.3-1] Proof header on its own line.
\makeatletter
\renewenvironment{proof}[1][\proofname]{\par
  \pushQED{\qed}\normalfont
  \topsep6\p@\@plus6\p@\relax
  \trivlist
  \item[\hskip\labelsep\itshape #1\@addpunct{.}]\mbox{}\par\nobreak\ignorespaces
}{\popQED\endtrivlist\@endpefalse}
\makeatother

\newcommand{\R}{\mathbb{R}}
\newcommand{\N}{\mathbb{N}}
\newcommand{\E}{\mathbb{E}}
\newcommand{\Prob}{\mathbb{P}}
\newcommand{\cF}{\mathcal{F}}
\newcommand{\cB}{\mathcal{B}}
\newcommand{\cP}{\mathcal{P}}
\newcommand{\Om}{\Omega}
\newcommand{\Leb}{\mathrm{Leb}}
\newcommand{\TV}{\mathrm{TV}}
\newcommand{\Normal}{\mathcal{N}}
\newcommand{\xobs}{\bm{x}_{\mathrm{obs}}}
\newcommand{\xhid}{\bm{x}_{\mathrm{hid}}}
\newcommand{\bth}{\bm{\theta}}
\newcommand{\push}{{}_{\#}}
\DeclareMathOperator{\intr}{int}
\newcommand{\norm}[1]{\lVert #1 \rVert}
\newcommand{\abs}[1]{\lvert #1 \rvert}

\title{\textbf{Theory Notes for A-DCVAE}\\[2pt]
\large Foundations, and Theorem~A: Uniform Ergodicity of the Continuous Pseudo-Gibbs Chain}
\author{Backstage reference note (internal), Rev.\ 3}
\date{\today}

\begin{document}
\maketitle

\begin{abstract}
Internal, fully rigorous note. We fix all notation, give self-contained definitions of every
measure-theoretic and Markov-chain object used, record the standard analytic facts we invoke (each
proved or precisely cited), state and prove the one substantive external tool (Doeblin's
minorization theorem, via an explicit coupling kernel and Ionescu--Tulcea), construct the A-DCVAE
inference kernel and verify it \emph{is} a Markov kernel, then prove \emph{Theorem~A}: under a
compact state space (from the boundary projection) and strictly positive Gaussian transition noise,
the continuous pseudo-Gibbs chain is uniformly ergodic. All numbered results are citable by later
notes (rate bounds; steering/consistency of the invariant law).
\end{abstract}

\tableofcontents

% ====================================================================
%  Secs. 1--9 below are VERBATIM from Rev.2 (per the Rev.3 constraint).
% ====================================================================
\section{Notation and standing conventions}
\label{sec:notation}

We work throughout on Borel subsets of Euclidean space. The abstract results below need no
topological assumption; the only place topology enters is the trivial observation, in the proof of
Theorem~\ref{thm:A}, that the diagonal of $\Om\subseteq\R^d$ is closed and hence lies in
$\cF\otimes\cF$.

\begin{notation}[Spaces and measures]\leavevmode
\begin{itemize}
  \item $\R^d$ carries the Euclidean norm $\norm{\cdot}$ and Borel $\sigma$-algebra $\cB(\R^d)$.
        For $S\subseteq\R^d$, $\intr S$ is its interior, $\partial S$ its boundary.
  \item $\Om$ is a fixed \emph{state space}, a Borel subset of some $\R^d$; $\cF=\cB(\Om)$.
  \item $\Leb$ is Lebesgue measure. $\cP(\Om)$ is the set of probability measures on $(\Om,\cF)$;
        $\delta_x$ is the Dirac mass at $x$.
  \item For measurable $f:\Om\to\Om'$ and $\mu\in\cP(\Om)$, the pushforward $f\push\mu\in\cP(\Om')$
        is $(f\push\mu)(A)=\mu(f^{-1}(A))$.
  \item $\Normal(\,\cdot\,;\bm m,\Sigma)$ is the Gaussian density on $\R^k$ with mean $\bm m$,
        covariance $\Sigma\succ0$; $\Normal(\bm m,\Sigma)$ the corresponding law; $I_k$ the identity.
\end{itemize}
\end{notation}

\begin{notation}[Model objects]\leavevmode
The sparse observation $\xobs$ is \emph{fixed} throughout (frozen at inference) and suppressed.
\begin{itemize}
  \item $\xhid\in\R^{d_H}$ hidden-state trajectory; $\bth\in\R^{d_\Theta}$ parameter; joint state
        $\bm z=(\xhid,\bth)$.
  \item Decoders $g_\phi:\R^{d_\Theta}\times\R^{m_H}\to\R^{d_H}$ and
        $g_\psi:\R^{d_H}\times\R^{m_P}\to\R^{d_\Theta}$ (with $\xobs$ fixed).
  \item Latent priors $\bm z_H\sim\Normal(\bm0,I_{m_H})$, $\bm z_P\sim\Normal(\bm0,I_{m_P})$;
        injection scales $\sigma_H,\sigma_P>0$.
  \item Projections $\mathcal B_H:\R^{d_H}\to\Om_H$, $\mathcal B_\Theta:\R^{d_\Theta}\to\Om_\Theta$
        (Def.~\ref{def:proj}).
\end{itemize}
\end{notation}

% ====================================================================
\section{Basic definitions}
\label{sec:defs}

\begin{definition}[Markov transition kernel]
\label{def:kernel}
A \emph{Markov kernel} on $(\Om,\cF)$ is $T:\Om\times\cF\to[0,1]$ such that (i) $T(x,\cdot)\in\cP(\Om)$
for each $x$, and (ii) $x\mapsto T(x,A)$ is measurable for each $A\in\cF$. It acts by
$(\mu T)(A)=\int T(x,A)\,\mu(dx)$ and $(Tf)(x)=\int f(y)\,T(x,dy)$, with the $n$-step kernel given by
Chapman--Kolmogorov $T^{n+1}(x,A)=\int T^n(y,A)\,T(x,dy)$, and, by an immediate induction, the
semigroup identity $T^{n+m}(x,A)=\int T^n(y,A)\,T^m(x,dy)$ for all $n,m\ge0$ (with $T^0=\mathrm{id}$).
\end{definition}

\begin{definition}[Invariant measure]
\label{def:invariant}
$\pi\in\cP(\Om)$ is \emph{invariant} for $T$ if $\pi T=\pi$.
\end{definition}

\begin{definition}[Total variation distance]
\label{def:tv}
For $\mu,\nu\in\cP(\Om)$, $\ \norm{\mu-\nu}_\TV:=\sup_{A\in\cF}\abs{\mu(A)-\nu(A)}\in[0,1]$.
Write $\lambda=\mu-\nu$ with Jordan decomposition $\lambda=\lambda^+-\lambda^-$ and Hahn positive
set $P$ (so $\lambda^+=\lambda(\cdot\cap P)$, $\lambda^-=-\lambda(\cdot\cap P^c)$). Then
$\norm{\mu-\nu}_\TV=\lambda^+(\Om)$: the supremum is attained at $A=P$ giving
$\lambda(P)=\lambda^+(\Om)$, and no $A$ exceeds it since $\lambda(A)\le\lambda^+(A)\le\lambda^+(\Om)$.
As $\lambda(\Om)=0$ forces $\lambda^+(\Om)=\lambda^-(\Om)$, also
$\norm{\mu-\nu}_\TV=\tfrac12\big(\lambda^++\lambda^-\big)(\Om)=\tfrac12\abs{\mu-\nu}(\Om)$. (Some
authors use twice this value.)
\end{definition}

\begin{definition}[Coupling]
\label{def:coupling}
A \emph{coupling} of $\mu,\nu$ is a pair $(X,Y)$ of $\Om$-valued random variables on one probability
space with $\mathrm{law}(X)=\mu$, $\mathrm{law}(Y)=\nu$.
\end{definition}

\begin{definition}[Minorization / Doeblin condition]
\label{def:minor}
$T$ satisfies the \emph{minorization condition} if there exist $\varepsilon\in(0,1]$ and
$\nu\in\cP(\Om)$ with $T(x,A)\ge\varepsilon\,\nu(A)$ for all $x\in\Om,\,A\in\cF$.
\end{definition}

\begin{definition}[Uniform ergodicity]
\label{def:unifergodic}
$T$ is \emph{uniformly ergodic} if there exist $\pi\in\cP(\Om)$, $C<\infty$, $\rho\in(0,1)$ with
$\sup_{x\in\Om}\norm{T^n(x,\cdot)-\pi}_\TV\le C\rho^n$ for all $n$. The supremum over $x$ (rate
independent of the start) distinguishes \emph{uniform} from merely geometric ergodicity.
\end{definition}

\begin{definition}[Euclidean projection onto a convex body]
\label{def:proj}
For nonempty compact convex $K\subset\R^d$, $\mathcal B_K(y):=\arg\min_{u\in K}\norm{y-u}$.
By Fact~\ref{fact:proj} this is single-valued, $1$-Lipschitz, fixes $K$ pointwise, and maps
$\R^d\setminus K$ into $\partial K$.
\end{definition}

% ====================================================================
\section{Standard facts used}
\label{sec:facts}

We isolate the classical facts the proofs rely on, each proved or precisely cited, so the note is
self-contained.

\begin{fact}[Projection onto a convex body]
\label{fact:proj}
For nonempty closed convex $K\subset\R^d$: $\mathcal B_K$ is well defined and single-valued;
$\norm{\mathcal B_K(y)-\mathcal B_K(y')}\le\norm{y-y'}$; $\mathcal B_K(y)=y\iff y\in K$; and
$y\notin K\Rightarrow\mathcal B_K(y)\in\partial K$. Consequently, for any $A\subseteq\intr K$ one has
$\mathcal B_K^{-1}(A)=A$.
\end{fact}
\begin{proof}
Existence/uniqueness and nonexpansiveness are standard (strict convexity of $\norm{\cdot}^2$ on the
convex $K$; see e.g.\ \cite[Ch.~III]{HiriartLemarechal}). If $y\in K$ the minimizer is $y$. If
$y\notin K$, let $u^\ast=\mathcal B_K(y)$ and suppose $u^\ast\in\intr K$; then for the segment
$u_t:=u^\ast+t(y-u^\ast)$ we have (i) $u_t\in K$ for all small $t>0$ (openness of $\intr K$ around
$u^\ast$; convexity is not even needed), and (ii)
$\norm{y-u_t}=(1-t)\norm{y-u^\ast}<\norm{y-u^\ast}$, since $y\neq u^\ast$ (because $y\notin K\ni u^\ast$).
This contradicts minimality, so $u^\ast\in\partial K$. For the last claim:
if $y\in A\subseteq\intr K\subseteq K$ then $\mathcal B_K(y)=y\in A$; if $y\notin K$ then
$\mathcal B_K(y)\in\partial K$, disjoint from $\intr K\supseteq A$, so $y\notin\mathcal B_K^{-1}(A)$;
if $y\in K\setminus A$ then $\mathcal B_K(y)=y\notin A$. Hence $\mathcal B_K^{-1}(A)=A$.
\end{proof}

\begin{fact}[Null boundary of a convex body]
\label{fact:nullbdry}
If $K\subset\R^d$ is compact convex with $\intr K\neq\emptyset$, then $\Leb(\partial K)=0$; in
particular $K$ is Lebesgue--Jordan measurable and $\Leb(\intr K)=\Leb(K)\in(0,\infty)$.
\end{fact}
\begin{proof}
After a translation assume $\bm 0\in\intr K$, so $B(\bm0,r)\subseteq K$ for some $r>0$. For
$\epsilon\in(0,1)$ and $x\in(1-\epsilon)K$, $\ x+\epsilon B(\bm0,r)\subseteq(1-\epsilon)K+\epsilon K\subseteq K$
by convexity, so $x\in\intr K$; hence $(1-\epsilon)K\subseteq\intr K$ and thus
$\partial K\cap(1-\epsilon)K=\emptyset$, i.e.\ $\partial K\subseteq K\setminus(1-\epsilon)K$.
Therefore $\Leb(\partial K)\le\Leb(K)-\Leb((1-\epsilon)K)=\big(1-(1-\epsilon)^d\big)\Leb(K)\to0$ as
$\epsilon\to0$, so $\Leb(\partial K)=0$. Finiteness and positivity of $\Leb(K)$ hold since $K$ is
compact with nonempty interior.
\end{proof}

\begin{fact}[Interior of a product]
\label{fact:prodint}
For $A\subseteq\R^m$, $B\subseteq\R^n$, $\ \intr(A\times B)=\intr A\times\intr B$.
Hence $\Leb(\partial\Om)=0$ and $\intr\Om=\intr\Om_H\times\intr\Om_\Theta$ for
$\Om=\Om_H\times\Om_\Theta$, using Fact~\ref{fact:nullbdry}.
\end{fact}
\begin{proof}
$(\supseteq)$: $\intr A\times\intr B$ is open (product of opens) and contained in $A\times B$, hence
in $\intr(A\times B)$. $(\subseteq)$: if $(a,b)\in\intr(A\times B)$, a basic open box
$U\times V\subseteq A\times B$ contains $(a,b)$ with $U,V$ open; then $a\in U\subseteq A$ and
$b\in V\subseteq B$, so $a\in\intr A$, $b\in\intr B$. The measure statement then follows from
Fact~\ref{fact:nullbdry} and
$\partial(A\times B)\subseteq(\partial A\times \bar B)\cup(\bar A\times\partial B)$.
\end{proof}

\begin{lemma}[Coupling inequality]
\label{lem:coupling}
For any coupling $(X,Y)$ of $\mu,\nu$, $\ \norm{\mu-\nu}_\TV\le\Prob(X\neq Y)$.
\end{lemma}
\begin{proof}
For $A\in\cF$, $\mu(A)-\nu(A)=\Prob(X\in A)-\Prob(Y\in A)\le\Prob(X\in A,Y\notin A)\le\Prob(X\neq Y)$;
symmetrise and take $\sup_A$.
\end{proof}

\begin{lemma}[The measure action is a TV-contraction]
\label{lem:nonexp}
For any Markov kernel $T$ and $\mu,\nu\in\cP(\Om)$, $\ \norm{\mu T-\nu T}_\TV\le\norm{\mu-\nu}_\TV$.
In particular $\mu\mapsto\mu T$ is $\norm{\cdot}_\TV$-continuous.
\end{lemma}
\begin{proof}
Let $\lambda=\mu-\nu$ with Jordan decomposition $\lambda=\lambda^+-\lambda^-$,
$\lambda^\pm(\Om)=\norm{\mu-\nu}_\TV$. For any $A$, since $0\le T(\cdot,A)\le1$,
$(\mu T-\nu T)(A)=\int T(x,A)\,\lambda(dx)\le\int T(x,A)\,\lambda^+(dx)\le\lambda^+(\Om)=\norm{\mu-\nu}_\TV$,
and symmetrically $\ge-\norm{\mu-\nu}_\TV$. Take $\sup_A$.
\end{proof}

\begin{lemma}[Averaging bound]
\label{lem:avg}
Let $P$ be a Markov kernel and $\eta\in\cP(\Om)$ a fixed measure with
$\norm{P(y,\cdot)-\eta}_\TV\le\beta$ for all $y\in\Om$. Then, for every $\rho\in\cP(\Om)$,
$\ \norm{\rho P-\eta}_\TV\le\beta$, where $\rho P:=\int P(y,\cdot)\,\rho(dy)$.
\end{lemma}
\begin{proof}
For each $A\in\cF$, using that $y\mapsto P(y,A)$ is measurable (kernel property) and
$\eta(A)=\int\eta(A)\,\rho(dy)$,
\[
  \abs{(\rho P)(A)-\eta(A)}=\biggl|\int\big(P(y,A)-\eta(A)\big)\rho(dy)\biggr|
  \le\int\abs{P(y,A)-\eta(A)}\,\rho(dy)\le\int\beta\,\rho(dy)=\beta .
\]
Take $\sup_A$. (Only pointwise bounds and kernel measurability are used; no measurability of
$y\mapsto\norm{\cdot}_\TV$ is needed.)
\end{proof}

\begin{fact}[Section integration against a kernel is measurable]
\label{fact:compose}
Let $Q$ be a Markov kernel from $(\mathsf S,\mathcal S)$ to $(\mathsf S',\mathcal S')$ and
$h:\mathsf S\times\mathsf S'\to\R$ bounded and $\mathcal S\otimes\mathcal S'$-measurable. Then
$s\mapsto\int h(s,t)\,Q(dt\mid s)$ is $\mathcal S$-measurable. In particular, for
$A\in\mathcal S\otimes\mathcal S'$, $\ s\mapsto Q\big(\{t:(s,t)\in A\}\mid s\big)$ is measurable; and
$(x,y)\mapsto\big(R(x,\cdot)\otimes R(y,\cdot)\big)(C)$ is jointly measurable for any Markov kernel
$R$ and $C\in\cF\otimes\cF$.
\end{fact}
\begin{proof}
Functional monotone class. Let $\mathcal H=\{h\text{ bdd meas.}:s\mapsto\int h(s,t)Q(dt\mid s)\text{ is
measurable}\}$. It is a vector space containing the constant $\mathbf 1=\mathbf 1_{\mathsf S\times\mathsf S'}$
(as $Q(\cdot\mid s)$ is a probability measure) and every rectangle indicator $\mathbf 1_{A\times B}$
(since $s\mapsto\mathbf 1_A(s)\,Q(B\mid s)$ is measurable by the kernel property), and is closed under
bounded monotone limits (monotone convergence). Rectangles form a $\pi$-system generating
$\mathcal S\otimes\mathcal S'$, so by the functional monotone class theorem $\mathcal H$ contains all
bounded $\mathcal S\otimes\mathcal S'$-measurable functions. Taking $h=\mathbf 1_A$ gives the first
special case. For the product-kernel case we apply the result twice, each time regarding a kernel as
one that ignores an extra coordinate: first, viewing $R(y,\cdot)$ as the kernel
$\tilde Q\big((u,y),\cdot\big):=R(y,\cdot)$ from $\mathsf S\times\mathsf S$, the map
$g(u,y):=R(y,C_u)=\int\mathbf 1_C(u,t)\,\tilde Q((u,y),dt)$ is jointly measurable (where
$C_u=\{t:(u,t)\in C\}$); second, viewing $R(x,\cdot)$ as $\hat Q\big((x,y),\cdot\big):=R(x,\cdot)$,
the map $(x,y)\mapsto\int g(u,y)\,\hat Q((x,y),du)=\big(R(x,\cdot)\otimes R(y,\cdot)\big)(C)$ is
measurable. (Cf.\ \cite[Ch.~1]{Kallenberg}.)
\end{proof}

\begin{corollary}[Powers are kernels]
\label{cor:Tn}
For every $n\ge0$, $T^n$ is a Markov kernel; in particular $y\mapsto T^n(y,A)$ is measurable for
each $A\in\cF$.
\end{corollary}
\begin{proof}
Induction on $n$ ($T^0=\mathrm{id}$, $T^1=T$). If $T^{n-1}$ is a kernel then $h(x,y):=T^{n-1}(y,A)$
is bounded measurable, so $x\mapsto T^n(x,A)=\int T^{n-1}(y,A)\,T(x,dy)$ is measurable by
Fact~\ref{fact:compose}, and $T^n(x,\cdot)$ is a probability measure.
\end{proof}

\begin{fact}[TV-completeness of $\cP(\Om)$]
\label{fact:complete}
$(\cP(\Om),\norm{\cdot}_\TV)$ is a complete metric space: a $\norm{\cdot}_\TV$-Cauchy sequence
$(\mu_n)$ converges in $\TV$ to some $\mu\in\cP(\Om)$.
\end{fact}
\begin{proof}
Let $(\mu_n)$ be $\TV$-Cauchy and dominate by $m:=\sum_{n\ge1}2^{-n}\mu_n$, which is a probability
measure (countable additivity of the series of measures by Tonelli, and $m(\Om)=\sum2^{-n}=1$). Then
$\mu_n\ll m$, so by Radon--Nikodym $f_n:=d\mu_n/dm\ge0$ with $\int f_n\,dm=1$. Since
$\mu_n-\mu_k=(f_n-f_k)\,m$, the signed measure has total variation
$\abs{\mu_n-\mu_k}(\Om)=\int\abs{f_n-f_k}\,dm$, so by Def.~\ref{def:tv}
$\norm{\mu_n-\mu_k}_\TV=\tfrac12\int\abs{f_n-f_k}\,dm=\tfrac12\norm{f_n-f_k}_{L^1(m)}$, and $(f_n)$ is
$L^1(m)$-Cauchy. By completeness of $L^1(m)$ (Riesz--Fischer, see e.g.\ \cite{Folland}), $f_n\to f$
in $L^1(m)$ for some $f\ge0$ a.e.\ with $\int f\,dm=\lim_n\int f_n\,dm=1$. Put $\mu:=f\,m\in\cP(\Om)$;
then $\norm{\mu_n-\mu}_\TV=\tfrac12\norm{f_n-f}_{L^1(m)}\to0$.
\end{proof}

\begin{fact}[Ionescu--Tulcea]
\label{fact:it}
Given \emph{any} measurable space $(\mathsf S,\mathcal S)$, an initial law, and a Markov kernel
$\bar T$ on $\mathsf S$, there exists a probability measure on $(\mathsf S^{\N},\mathcal S^{\otimes\N})$
under which the coordinate process is a Markov chain with transition $\bar T$
\cite[Ch.~6]{Kallenberg}. Unlike Kolmogorov's extension theorem, no topological (standard Borel)
assumption is needed.
\end{fact}

% ====================================================================
\section{Analytic tool: Doeblin's theorem}
\label{sec:tool}

\begin{theorem}[Doeblin; see {\cite[Ch.~16]{MeynTweedie}}]
\label{thm:doeblin}
Let $T$ be a Markov kernel on a measurable space $(\Om,\cF)$ whose diagonal
$\{(x,x):x\in\Om\}\in\cF\otimes\cF$, and suppose $T$ satisfies the minorization condition
(Def.~\ref{def:minor}) with $(\varepsilon,\nu)$. Then $T$ has a unique invariant
$\pi\in\cP(\Om)$, and $\ \sup_{x\in\Om}\norm{T^n(x,\cdot)-\pi}_\TV\le(1-\varepsilon)^n$ for all
$n\ge0$; consequently $\norm{\mu_0T^n-\pi}_\TV\le(1-\varepsilon)^n$ for every $\mu_0\in\cP(\Om)$.
\end{theorem}

\begin{proof}
The case $n=0$ holds trivially under the convention $(1-\varepsilon)^0=1$, since $\norm{\cdot}_\TV\le1$
always (Def.~\ref{def:tv}); we treat $n\ge1$ throughout.

\emph{Degenerate case $\varepsilon=1$.} Then $T(x,A)\ge\nu(A)$ for all $A$, and since $T(x,\cdot),\nu$
are both probability measures this forces $T(x,\cdot)=\nu$ for every $x$; hence
$T^n(x,\cdot)=\nu=:\pi$ for all $n\ge1$, which is invariant, unique, and satisfies the bound
$(1-1)^n=0$. Assume henceforth $\varepsilon\in(0,1)$.

\emph{Splitting.} $R(x,A):=\dfrac{T(x,A)-\varepsilon\nu(A)}{1-\varepsilon}$ is, for each $x$, a
probability measure (nonnegative by minorization, total mass $1$) and is measurable in $x$; thus $R$
is a Markov kernel and
\begin{equation}\label{eq:split}
  T(x,\cdot)=\varepsilon\,\nu(\cdot)+(1-\varepsilon)\,R(x,\cdot).
\end{equation}

\emph{Coupling kernel.} Define $\bar T$ on the product space $(\Om\times\Om,\cF\otimes\cF)$ by
\[
  \bar T\big((x,y),\cdot\big)=
  \begin{cases}
    \varepsilon\,(\Delta\push\nu)+(1-\varepsilon)\,\big(R(x,\cdot)\otimes R(y,\cdot)\big), & x\neq y,\\[2pt]
    \Delta\push T(x,\cdot), & x=y,
  \end{cases}
\]
where $\Delta(z)=(z,z)$. The diagonal $D:=\{x=y\}$ lies in $\cF\otimes\cF$ by the standing hypothesis
on $(\Om,\cF)$, so the two-branch definition is legitimate. The off-diagonal branch
$(x,y)\mapsto R(x,\cdot)\otimes R(y,\cdot)$ is jointly measurable by Fact~\ref{fact:compose}; the
diagonal branch $x\mapsto\Delta\push T(x,\cdot)$ is measurable because, for $C\in\cF\otimes\cF$,
$(\Delta\push T(x,\cdot))(C)=T(x,\Delta^{-1}(C))$ with $\Delta^{-1}(C)\in\cF$ ($\Delta$ measurable),
and $x\mapsto T(x,\Delta^{-1}(C))$ is measurable (kernel property). Hence $\bar T$ is a Markov kernel
on $(\Om\times\Om,\cF\otimes\cF)$. By Ionescu--Tulcea (Fact~\ref{fact:it}, valid on any measurable
space) there is, for each start $(x,y)$, a Markov chain $(X_n,Y_n)_{n\ge0}$ with transition $\bar T$
and $(X_0,Y_0)=(x,y)$; let $(\mathcal F_n)$ be its natural filtration. (It is precisely here that the
measurable-diagonal hypothesis is used: it makes the event $\{X_{n+1}=Y_{n+1}\}\in\cF\otimes\cF$, so
the conditional probability below is well defined.)

\emph{Marginals and coalescence.} By \eqref{eq:split}, projecting $\bar T$ to the first coordinate
yields, in both branches, exactly $T(x,\cdot)$; likewise the second coordinate yields $T(y,\cdot)$.
By induction $\mathrm{law}(X_n)=T^n(x,\cdot)$ and $\mathrm{law}(Y_n)=T^n(y,\cdot)$. Whenever
$X_n=Y_n$ the $x=y$ branch keeps the coordinates equal forever. Let $\tau=\inf\{n:X_n=Y_n\}$. On the
event $\{X_n\neq Y_n\}$ the $\varepsilon$-branch sends both coordinates to a common $Z\sim\nu$, so
$\Prob(X_{n+1}=Y_{n+1}\mid\mathcal F_n)\ge\varepsilon$; hence
\[
  \Prob(\tau>n+1)=\E\big[\mathbf 1_{\{\tau>n\}}\,\Prob(X_{n+1}\neq Y_{n+1}\mid\mathcal F_n)\big]
  \le(1-\varepsilon)\,\Prob(\tau>n),
\]
and by induction $\Prob(\tau>n)\le(1-\varepsilon)^n$. Since $X_n=Y_n$ on $\{\tau\le n\}$, the
coupling inequality (Lemma~\ref{lem:coupling}) gives
\begin{equation}\label{eq:contract}
  \norm{T^n(x,\cdot)-T^n(y,\cdot)}_\TV\le\Prob(X_n\neq Y_n)\le\Prob(\tau>n)\le(1-\varepsilon)^n .
\end{equation}

Note \eqref{eq:contract} says exactly $\norm{T^n(y,\cdot)-\eta}_\TV\le(1-\varepsilon)^n$ for the fixed
$\eta=T^n(x,\cdot)$ and all $y$; this is the hypothesis of the averaging bound
(Lemma~\ref{lem:avg}) with $P=T^n$.

\emph{Existence.} Fix $x$. Apply Lemma~\ref{lem:avg} with $\eta=T^n(x,\cdot)$, $P=T^n$,
$\rho=T^m(x,\cdot)$: since $\rho P=T^m(x,\cdot)T^n=T^{n+m}(x,\cdot)$ by the semigroup identity
(Def.~\ref{def:kernel}), this yields $\norm{T^{n+m}(x,\cdot)-T^n(x,\cdot)}_\TV\le(1-\varepsilon)^n$
for all $m$. So $(T^n(x,\cdot))_n$ is $\TV$-Cauchy; by Fact~\ref{fact:complete} it converges to some
$\pi\in\cP(\Om)$. Applying the $\TV$-continuous map $\mu\mapsto\mu T$ (Lemma~\ref{lem:nonexp}) to
$T^{n}(x,\cdot)\to\pi$ gives $T^{n+1}(x,\cdot)\to\pi T$; but the left side also $\to\pi$, so
$\pi T=\pi$.

\emph{Uniqueness and the bounds.} Let $\pi'$ be invariant. Apply Lemma~\ref{lem:avg} with
$\eta=T^n(x,\cdot)$, $P=T^n$, $\rho=\pi'$: since $\rho P=\pi'T^n=\pi'$ by invariance, we get
$\norm{\pi'-T^n(x,\cdot)}_\TV\le(1-\varepsilon)^n$; letting $n\to\infty$ forces $\pi'=\pi$, and the
inequality (for every $x$) is precisely the $\sup_x$ bound $\sup_x\norm{T^n(x,\cdot)-\pi}_\TV\le(1-\varepsilon)^n$.
Finally, that $\sup_x$ bound reads $\norm{T^n(y,\cdot)-\pi}_\TV\le(1-\varepsilon)^n$ for all $y$; a last
application of Lemma~\ref{lem:avg} with $\eta=\pi$, $P=T^n$, $\rho=\mu_0$ gives
$\norm{\mu_0T^n-\pi}_\TV\le(1-\varepsilon)^n$.
\end{proof}

% ====================================================================
\section{The A-DCVAE inference kernel}
\label{sec:kernel}

\begin{definition}[Conditional stochastic maps]
\label{def:condmaps}
Define the post-projection conditional laws
\[
  Q_H(\cdot\mid\bth):=(\mathcal B_H)\push\ \mathrm{law}\big(g_\phi(\bth,\bm z_H)+\sigma_H\bm\xi_H\big),
  \quad \bm\xi_H\sim\Normal(\bm0,I_{d_H}),
\]
\[
  Q_P(\cdot\mid\xhid'):=(\mathcal B_\Theta)\push\ \mathrm{law}\big(g_\psi(\xhid',\bm z_P)+\sigma_P\bm\xi_P\big),
  \quad \bm\xi_P\sim\Normal(\bm0,I_{d_\Theta}),
\]
with $\bm z_H,\bm\xi_H,\bm z_P,\bm\xi_P$ mutually independent. The pre-projection law of the hidden
step has Lebesgue density on $\R^{d_H}$
\begin{equation}\label{eq:qH-density}
  q_H(\bm u\mid\bth)=\int_{\R^{m_H}}\Normal\big(\bm u;\,g_\phi(\bth,\bm z_H),\,\sigma_H^2 I_{d_H}\big)\,
  \Normal(\bm z_H;\bm0,I_{m_H})\,d\bm z_H,
\end{equation}
an infinite Gaussian mixture; $q_P(\cdot\mid\xhid')$ is analogous.
\end{definition}

\begin{definition}[The inference kernel $T$]
\label{def:T}
On $\Om:=\Om_H\times\Om_\Theta$ with state $\bm z=(\xhid,\bth)$, one sweep draws
$\xhid'\sim Q_H(\cdot\mid\bth)$ then $\bth'\sim Q_P(\cdot\mid\xhid')$, giving
\[
  T\big((\xhid,\bth),\,A\big)=\int_{\Om_H} Q_P\big(A_{\bm u'}\mid\bm u'\big)\,Q_H(d\bm u'\mid\bth),
  \qquad A_{\bm u'}:=\{\bm t':(\bm u',\bm t')\in A\}.
\]
The transition depends on the current state only through $\bth$.
\end{definition}

\begin{lemma}[$T$ is a Markov kernel]
\label{lem:measurable}
Under Assumptions~\ref{ass:compact}--\ref{ass:cont}, $Q_H,Q_P$ are Markov kernels (between the
respective spaces) and $T$ of Def.~\ref{def:T} is a Markov kernel on $(\Om,\cF)$.
\end{lemma}
\begin{proof}
\emph{Joint measurability of the density.} In \eqref{eq:qH-density} the integrand is continuous in
$(\bm u,\bth,\bm z_H)$ (Gaussian density and $g_\phi$ continuous) and dominated by the integrable
$(2\pi\sigma_H^2)^{-d_H/2}\Normal(\bm z_H;\bm0,I_{m_H})$; by dominated convergence
$(\bm u,\bth)\mapsto q_H(\bm u\mid\bth)$ is continuous, hence jointly Borel measurable. The same
formula gives the pre-projection density on all of $\R^{d_H}$.

\emph{$Q_H$ is a kernel.} For each $\bth$, $Q_H(\cdot\mid\bth)$ is a probability measure (pushforward
of one under the continuous $\mathcal B_H$). For measurability in $\bth$, write via the pushforward
representation
\[
  Q_H(A\mid\bth)=\int_{\R^{d_H}}\mathbf 1_A\big(\mathcal B_H(\bm y)\big)\,q_H(\bm y\mid\bth)\,d\bm y ,
\]
whose integrand is jointly measurable in $(\bm y,\bth)$ ($\mathcal B_H$ continuous, $A$ Borel, $q_H$
jointly measurable); by Tonelli $\bth\mapsto Q_H(A\mid\bth)$ is measurable. Thus $Q_H$ is a Markov
kernel from $\Om_\Theta$ to $\Om_H$; symmetrically $Q_P$ is a kernel from $\Om_H$ to $\Om_\Theta$.

\emph{$T$ is a kernel.} For fixed $\bth$, $A\mapsto T((\xhid,\bth),A)$ is a probability measure
(composition of the kernels $Q_H,Q_P$). For measurability of $\bm z\mapsto T(\bm z,A)$, note
$\bm u'\mapsto Q_P(A_{\bm u'}\mid\bm u')=Q_P(\{\bm t':(\bm u',\bm t')\in A\}\mid\bm u')$ is measurable
by Fact~\ref{fact:compose}; integrating this against the kernel $Q_H(\cdot\mid\bth)$ (again
Fact~\ref{fact:compose}) preserves measurability in $\bth$, hence in $\bm z$.
\end{proof}

% ====================================================================
\section{Standing assumptions}
\label{sec:assump}

\begin{assumption}[Compact convex state space via projection]
\label{ass:compact}
$\Om_H\subset\R^{d_H}$, $\Om_\Theta\subset\R^{d_\Theta}$ are nonempty, compact, convex with
nonempty interior, and $\mathcal B_H,\mathcal B_\Theta$ are the Euclidean projections onto them.
Thus $\Om=\Om_H\times\Om_\Theta$ is compact, standard Borel, with $\Leb(\Om)\in(0,\infty)$ and
$\Leb(\partial\Om)=0$ (Facts~\ref{fact:nullbdry}--\ref{fact:prodint}).
\end{assumption}

\begin{assumption}[Continuity and nondegenerate noise]
\label{ass:cont}
For the fixed $\xobs$, $g_\phi$ and $g_\psi$ are continuous on their domains and
$\sigma_H,\sigma_P>0$.
\end{assumption}

\noindent Both hold \emph{by construction}: projections give Assumption~\ref{ass:compact};
continuous-activation networks with positive sampling noise give Assumption~\ref{ass:cont}. This is
the sense in which ergodicity is \emph{proved} from what the algorithm controls.

% ====================================================================
\section{Minorization lemma}
\label{sec:lemma}

\begin{lemma}[Uniform minorization of $T$]
\label{lem:minor}
Under Assumptions~\ref{ass:compact}--\ref{ass:cont}, $T$ satisfies the minorization condition with
$\nu(\cdot)=\Leb(\cdot\cap\intr\Om)/\Leb(\Om)$ and some $\varepsilon\in(0,1]$.
\end{lemma}
\begin{proof}
\emph{Step 1 (uniform density lower bound).} Fix $R>0$ with $p_R:=\Prob(\norm{\bm z_H}\le R)>0$. As
$\Om_\Theta\times\{\norm{\bm z_H}\le R\}$ is compact and $g_\phi$ continuous,
$M_H:=\sup\{\norm{g_\phi(\bth,\bm z_H)}:\bth\in\Om_\Theta,\norm{\bm z_H}\le R\}<\infty$; let
$M_H':=\sup_{\bm u\in\Om_H}\norm{\bm u}<\infty$. For $\bm u\in\Om_H,\bth\in\Om_\Theta,\norm{\bm z_H}\le R$,
\[
  \Normal(\bm u;g_\phi(\bth,\bm z_H),\sigma_H^2 I_{d_H})
  \ge(2\pi\sigma_H^2)^{-d_H/2}\exp\!\Big(-\tfrac{(M_H+M_H')^2}{2\sigma_H^2}\Big)=:c_H>0 .
\]
Restricting \eqref{eq:qH-density} to $\{\norm{\bm z_H}\le R\}$,
\begin{equation}\label{eq:deltaH}
  q_H(\bm u\mid\bth)\ge c_H\,p_R=:\delta_H>0\quad\text{uniformly over }\bm u\in\Om_H,\ \bth\in\Om_\Theta,
\end{equation}
and analogously $q_P(\bm t\mid\bm u')\ge\delta_P>0$ uniformly.

\emph{Step 2 (projection preserves the interior bound).} By Fact~\ref{fact:proj}, for
$A\subseteq\intr\Om_H$ we have $\mathcal B_H^{-1}(A)=A$, so
$Q_H(A\mid\bth)=\int_{\mathcal B_H^{-1}(A)}q_H(\bm y\mid\bth)\,d\bm y=\int_A q_H(\bm y\mid\bth)\,d\bm y\ge\delta_H\,\Leb(A)$
by \eqref{eq:deltaH}; likewise $Q_P(B\mid\bm u')\ge\delta_P\,\Leb(B)$ for $B\subseteq\intr\Om_\Theta$.
Equivalently, $Q_H(\cdot\mid\bth)-\delta_H\,\Leb(\cdot\cap\intr\Om_H)$ and
$Q_P(\cdot\mid\bm u')-\delta_P\,\Leb(\cdot\cap\intr\Om_\Theta)$ are nonnegative measures, so the
bounds pass to integrals of nonnegative functions.

\emph{Step 3 (composition).} For any $A\in\cF$ and state $(\xhid,\bth)$, using
$\intr\Om=\intr\Om_H\times\intr\Om_\Theta$ (Fact~\ref{fact:prodint}), the nonnegative-measure
domination from Step 2, and Tonelli,
\[
  T\big((\xhid,\bth),A\big)=\int_{\Om_H}Q_P(A_{\bm u'}\mid\bm u')\,Q_H(d\bm u'\mid\bth)
  \ge\int_{\intr\Om_H}\!\!\delta_P\,\Leb\big(A_{\bm u'}\cap\intr\Om_\Theta\big)\,\delta_H\,d\bm u'
  =\delta_H\delta_P\,\Leb(A\cap\intr\Om).
\]
Set $\varepsilon:=\delta_H\delta_P\,\Leb(\Om)$ and $\nu(\cdot)=\Leb(\cdot\cap\intr\Om)/\Leb(\Om)$, so
$T((\xhid,\bth),A)\ge\varepsilon\,\nu(A)$ for all states and $A$. Taking $A=\Om$ gives
$1=T((\xhid,\bth),\Om)\ge\varepsilon\,\nu(\Om)=\varepsilon$ (as $\nu(\Om)=1$ by
$\Leb(\partial\Om)=0$), so $\varepsilon\in(0,1]$.
\end{proof}

% ====================================================================
\section{Theorem A: uniform ergodicity of the continuous Pseudo-Gibbs chain}
\label{sec:thmA}

\begin{theorem}[Uniform ergodicity / global convergence]
\label{thm:A}
Under Assumptions~\ref{ass:compact}--\ref{ass:cont}, the inference kernel $T$ (Def.~\ref{def:T}) has
a unique invariant $\pi^\ast\in\cP(\Om)$, and for every initial $\mu_0\in\cP(\Om)$,
\[
  \norm{\mu_0\,T^n-\pi^\ast}_\TV\le(1-\varepsilon)^n,\qquad n\in\N,
\]
with $\varepsilon=\delta_H\delta_P\Leb(\Om)\in(0,1]$ from Lemma~\ref{lem:minor}. The rate is uniform
over all initializations; equivalently $\mu\mapsto\mu T$ on $\cP(\Om)$ has a globally attracting
fixed point $\pi^\ast$.
\end{theorem}
\begin{proof}
By Lemma~\ref{lem:measurable} $T$ is a Markov kernel on $(\Om,\cF)$. Since $\Om\subseteq\R^d$ is a
separable metric space, its diagonal $\{(x,x):x\in\Om\}$ is closed, hence Borel, and therefore lies
in $\cF\otimes\cF$; thus the standing hypothesis of Theorem~\ref{thm:doeblin} holds. By
Lemma~\ref{lem:minor}, $T$ satisfies that theorem's minorization condition with $(\varepsilon,\nu)$.
The conclusion is Theorem~\ref{thm:doeblin}'s.
\end{proof}

% ====================================================================
\section{Remarks}
\label{sec:remarks}

\begin{remark}[Continuous extension of Heckerman's pseudo-Gibbs]
\label{rem:heckerman}
For \emph{finite} state spaces a strictly positive transition matrix gives a unique stationary law
and geometric convergence at once by Perron--Frobenius; this is the regime of dependency networks.
Theorem~\ref{thm:A} is the continuous counterpart, where positivity alone is insufficient and the
minorization of Lemma~\ref{lem:minor} supplies the missing ingredient.
\end{remark}

\begin{remark}[Role of the noise; what deterministic collapse does and does not mean]
\label{rem:noise}
Doeblin's condition is \emph{sufficient}, not necessary, and the minorizing measure $\nu$ need not
be absolutely continuous: e.g.\ a constant decoder gives $T(x,\cdot)=\delta_c$, which satisfies
minorization with $\nu=\delta_c,\varepsilon=1$ and is (trivially) uniformly ergodic. Thus removing
the noise ($\sigma_H=\sigma_P=0$) does not by itself falsify Theorem~\ref{thm:A}; what it does is
break Assumption~\ref{ass:cont} and the proof route of Lemma~\ref{lem:minor} (each conditional law
becomes a Dirac mass, with no Lebesgue-density lower bound). Moreover, for a deterministic
contraction one has $\delta_{x_n}\to\delta_{x^\ast}$ \emph{weakly} but in general \emph{not} in
total variation, so TV-uniform ergodicity generally fails in the deterministic case. The injected
Gaussian noise is precisely what makes our Lemma~\ref{lem:minor} (and hence the TV guarantee)
available.
\end{remark}

\begin{remark}[What is, and is not, controlled]
\label{rem:scope}
Theorem~\ref{thm:A} concerns convergence of the \emph{distribution} $\mu_0T^n$, not of any sample
path. The constant $\varepsilon=\delta_H\delta_P\Leb(\Om)$ is typically minuscule: from
\eqref{eq:deltaH}, $\delta_H=(2\pi\sigma_H^2)^{-d_H/2}p_R\,\exp(-(M_H+M_H')^2/2\sigma_H^2)$, whose
prefactor $(2\pi\sigma_H^2)^{-d_H/2}$ and exponent are both strongly dimension- and
$\sigma$-dependent (the curse of dimensionality enters here). Hence $(1-\varepsilon)^n$ certifies
convergence but is a loose rate; dimension-robust rates require the complementary Wasserstein
contraction under condition-noise Jacobian control (companion rate note). We therefore avoid writing
$\delta_H\gtrsim\exp(\cdots)$ without flagging the hidden prefactor.
\end{remark}

\begin{remark}[Forward pointers]
\label{rem:forward}
Theorem~\ref{thm:A} settles existence and reachability of $\pi^\ast$. Separate notes treat (i) the
\emph{rate} via $2$-Wasserstein contraction from Jacobian control, and (ii) \emph{steering /
consistency}: how far $\pi^\ast$ lies from the target joint posterior when $Q_H,Q_P$ are mutually
incompatible, and how target-denoising aligns $\pi^\ast$ with the data manifold. The objects fixed
here---the composition-measurability Fact~\ref{fact:compose} and the powers Corollary~\ref{cor:Tn},
the averaging bound (Lemma~\ref{lem:avg}), the density bounds \eqref{eq:deltaH}, the measurability
Lemma~\ref{lem:measurable}, and Lemmas~\ref{lem:coupling}--\ref{lem:nonexp}---are reused there;
Fact~\ref{fact:compose} in particular recurs in any kernel manipulation.
\end{remark}

% ====================================================================
%  From here on (Secs. 10--13): Rev.3 edits, each marked [Rev.3-n].
% ====================================================================
\section{Wasserstein preliminaries}
\label{sec:wass}

The existence result (Theorem~\ref{thm:A}) lives in total variation. For the \emph{rate} and, later,
the stability/steering bound, TV is the wrong metric: an \emph{injective} deterministic map preserves
total variation (it maps two distinct Diracs to two distinct Diracs, still at TV-distance $1$), so it
cannot TV-contract, and more generally TV is blind to how close distinct points are. (Non-injective
maps such as a constant \emph{can} TV-contract, cf.\ Remark~\ref{rem:noise}; the point is only that TV
does not see the geometry that governs mixing.) We therefore switch to the quadratic Wasserstein
distance, which Lipschitz maps \emph{do} contract.

\begin{definition}[Quadratic Wasserstein distance]
\label{def:w2}
For $\mu,\nu\in\cP(\Om)$ on compact $\Om\subseteq\R^d$,
\[
  W_2(\mu,\nu):=\Big(\inf_{\gamma\in\Gamma(\mu,\nu)}\int_{\Om\times\Om}\norm{x-y}^2\,\gamma(dx,dy)\Big)^{1/2},
\]
where $\Gamma(\mu,\nu)$ is the set of couplings (Def.~\ref{def:coupling}).
\end{definition}

\begin{fact}[Basic properties on a compact domain]
\label{fact:w2basic}
On compact $\Om$, $W_2$ is a finite metric; $(\cP(\Om),W_2)$ is a complete (indeed compact) metric
space and $W_2$ metrizes weak convergence \cite[Ch.~6]{Villani}. For \emph{any} coupling $(X,Y)$ of
$\mu,\nu$, $\ W_2(\mu,\nu)\le(\E\norm{X-Y}^2)^{1/2}$ (the infimum is over couplings).
\end{fact}

\begin{fact}[Lipschitz maps contract $W_2$; marginalization is non-expansive]
\label{fact:w2lip}
If $f$ is $L$-Lipschitz then $W_2(f\push\mu,f\push\nu)\le L\,W_2(\mu,\nu)$ (push an optimal coupling
through $f\times f$). In particular the coordinate projection
$\mathrm{proj}_\Theta:\Om_H\times\Om_\Theta\to\Om_\Theta$ is $1$-Lipschitz, so
$W_2\big(\mathrm{proj}_\Theta\push\mu,\ \mathrm{proj}_\Theta\push\nu\big)\le W_2(\mu,\nu)$; this yields
the $\bth$-marginal form of any joint $W_2$ bound for free.
\end{fact}

\begin{fact}[Banach fixed point in $(\cP(\Om),W_2)$]
\label{fact:banachw2}
If a Markov kernel's measure action $\mu\mapsto\mu K$ satisfies $W_2(\mu K,\nu K)\le\kappa\,W_2(\mu,\nu)$
with $\kappa\in[0,1)$, then $K$ has a unique invariant $\nu^\ast$ and
$W_2(\mu_0K^n,\nu^\ast)\le\kappa^n W_2(\mu_0,\nu^\ast)\to0$ (Banach fixed point on the complete space
of Fact~\ref{fact:w2basic}).
\end{fact}

% ---------------------------------------------------------------------
% [Rev.3-2] NEW FACT. WHY: Villani Cor. 5.22 requires s -> (K(s,.),K'(s,.))
% to be Borel for the WEAK topology on P(Om), whereas Def. 2.1 gives only
% measurability w.r.t. the evaluation sigma-algebra sigma(mu -> mu(A)).
% On a Polish target these coincide; recording this once closes the gap
% and also makes mixtures of coupling kernels (used in Fact msel(iii),
% Lemma sweep, Lemma perturb, Lemma opdisc) well defined.
% ---------------------------------------------------------------------
\begin{fact}[Kernels are Borel maps into $(\cP(\Om),\mathrm{weak})$; mixtures]
\label{fact:kernmeas}
Let the target space $\Om$ be Polish (e.g.\ compact $\Om\subset\R^d$, or $\Om\times\Om$). On
$\cP(\Om)$ the evaluation $\sigma$-algebra $\sigma\big(\mu\mapsto\mu(A):A\in\cF\big)$ coincides with
the Borel $\sigma$-algebra of the weak topology \cite[Chs.~1 and 4]{Kallenberg}. Consequently:
(i) a map $s\mapsto\mu_s$ from $(\mathsf S,\mathcal S)$ is a Markov kernel \emph{iff} it is Borel
measurable into $\cP(\Om)$ carrying the weak topology; and
(ii) for any Markov kernel $s\mapsto\gamma_s$ into $\cP(\Om)$ and any $\rho\in\cP(\mathsf S)$, the
\emph{mixture} $\bar\gamma:=\int\gamma_s\,\rho(ds)$, defined by $\bar\gamma(C):=\int\gamma_s(C)\,\rho(ds)$,
is a well-defined element of $\cP(\Om)$ (the integrand is measurable by the kernel property;
countable additivity follows by monotone convergence).
\end{fact}

% ---------------------------------------------------------------------
% [Rev.3-3] Fact msel restated as (i)-(iii) with a full proof.
% WHY/HOW: (a) the previous "equivalently ... measurable" claimed an
% equivalence where only one direction is needed; measurability of
% s -> W_2 is now PROVED directly (continuity of W_2 in the weak topology
% on compact Om, composed with the Borel map of Fact kernmeas);
% (b) applicability of Villani Cor. 5.22 is now justified via Fact
% kernmeas rather than asserted; (c) the operation int gamma_s rho(ds)
% is a MIXTURE, not an application of the gluing lemma -- terminology
% corrected here and at all call sites; (d) the mixture is renamed
% bar-gamma to avoid clashing with Gamma(mu,nu) of Def. 10.1.
% ---------------------------------------------------------------------
\begin{fact}[Measurable selection of optimal couplings; averaged $W_2$ of kernels]
\label{fact:msel}
Let $K,K'$ be Markov kernels from $(\mathsf S,\mathcal S)$ to the compact $\Om$. Then:
\emph{(i)} $s\mapsto W_2\big(K(s,\cdot),K'(s,\cdot)\big)$ is measurable;
\emph{(ii)} there is a Markov kernel $s\mapsto\gamma_s$ from $(\mathsf S,\mathcal S)$ to
$\Om\times\Om$ such that, for \emph{every} $s$, $\gamma_s$ is an \emph{optimal} coupling of
$K(s,\cdot)$ and $K'(s,\cdot)$ \cite[Cor.~5.22]{Villani};
\emph{(iii)} for every $\rho\in\cP(\mathsf S)$, the mixture $\bar\gamma:=\int\gamma_s\,\rho(ds)$ is a
coupling of $\rho K$ and $\rho K'$, and hence
\[
  W_2(\rho K,\rho K')^2\ \le\ \int W_2\big(K(s,\cdot),K'(s,\cdot)\big)^2\,\rho(ds)
  \ \le\ \sup_{s}\,W_2\big(K(s,\cdot),K'(s,\cdot)\big)^2 .
\]
\end{fact}
\begin{proof}
\emph{(i).} By Fact~\ref{fact:kernmeas}(i), $s\mapsto\big(K(s,\cdot),K'(s,\cdot)\big)$ is Borel into
$\cP(\Om)\times\cP(\Om)$ with the weak topology. On compact $\Om$, $W_2$ metrizes that topology
(Fact~\ref{fact:w2basic}), so $W_2:\cP(\Om)\times\cP(\Om)\to\R$ is continuous; the composition is
measurable.

\emph{(ii).} By Fact~\ref{fact:kernmeas}(i) the hypothesis of \cite[Cor.~5.22]{Villani} (Borel
measurability of $s\mapsto(K(s,\cdot),K'(s,\cdot))$, target spaces Polish) is met; that corollary
returns a Borel map $s\mapsto\gamma_s\in\cP(\Om\times\Om)$ with each $\gamma_s$ an optimal transport
plan for the quadratic cost. By Fact~\ref{fact:kernmeas}(i) again (applied to the Polish
$\Om\times\Om$), this Borel map \emph{is} a Markov kernel.

\emph{(iii).} $\bar\gamma\in\cP(\Om\times\Om)$ is well defined by Fact~\ref{fact:kernmeas}(ii). Its
marginals: for Borel $A\subseteq\Om$,
$\bar\gamma(A\times\Om)=\int\gamma_s(A\times\Om)\,\rho(ds)=\int K(s,A)\,\rho(ds)=(\rho K)(A)$, and
symmetrically $\bar\gamma(\Om\times A)=(\rho K')(A)$; so $\bar\gamma\in\Gamma(\rho K,\rho K')$.
Therefore, using optimality of each $\gamma_s$ and Tonelli,
\[
  W_2(\rho K,\rho K')^2\le\int_{\Om\times\Om}\norm{x-y}^2\,\bar\gamma(dx,dy)
  =\int_{\mathsf S}\!\int_{\Om\times\Om}\norm{x-y}^2\,\gamma_s(dx,dy)\,\rho(ds)
  =\int W_2\big(K(s,\cdot),K'(s,\cdot)\big)^2\rho(ds).
\]
The last inequality of the statement is immediate.
\end{proof}

% ====================================================================
\section{Conditional Lipschitzness and the sweep contraction}
\label{sec:contraction}

% [Rev.3-10] Notation in this and the following sections unified to the
% bold convention of Secs. 1--8 (\bth, \bm u, \bm z_H, ...). Purely
% notational; no mathematical content changed by this unification.
We now quantify the mixing \emph{rate} of the same kernel $T$. Recall (Def.~\ref{def:condmaps}) the
conditional laws
$Q_H(\cdot\mid\bth)=(\mathcal B_H)\push\,\mathrm{law}(g_\phi(\bth,\bm z_H)+\sigma_H\bm\xi_H)$ and
$Q_P(\cdot\mid\bm u)=(\mathcal B_\Theta)\push\,\mathrm{law}(g_\psi(\bm u,\bm z_P)+\sigma_P\bm\xi_P)$,
where $\bm z_H,\bm z_P$ are the CVAE latents and $\sigma_H,\sigma_P>0$ the injection-noise scales.

\begin{assumption}[Conditional Lipschitz decoders --- (A-contraction)]
\label{ass:lip}
There exist $L_H,L_P<\infty$ with, uniformly in the latent code $\bm z$,
\[
  \sup_{\bm z}\ \mathrm{Lip}_{\bth\in\Om_\Theta}\big(g_\phi(\cdot,\bm z)\big)\le L_H,\qquad
  \sup_{\bm z}\ \mathrm{Lip}_{\bm u\in\Om_H}\big(g_\psi(\cdot,\bm z)\big)\le L_P ,
\]
i.e.\ each decoder mean is Lipschitz in its \emph{conditioning} variable, over the compact domains
$\Om_\Theta,\Om_H$ (these contain the supports used below, cf.\ Assumption~\ref{ass:suptrue}). This is exactly the
quantity that condition-noise injection (N1) regularizes, by the noise-equals-Tikhonov identity
\cite{Bishop}: training penalizes $\E\|\nabla_{\text{cond}}g\|_F^2$, promoting small $L_H,L_P$. We
posit this as an assumption (the penalty encourages but does not certify a bound) and
support it empirically (autocorrelation decay). \emph{Crucially, $L_H,L_P$ constrain the dependence on
the condition only --- not on the latent $\bm z$ (N3) nor on the injection noise $\sigma\bm\xi$ (N4).}
\end{assumption}

% [Rev.3-9] Quantifier domains made explicit (theta in Om_Theta, u in
% Om_H), matching the Lipschitz domains fixed in Assumption lip. All
% call sites already live in these ranges.
\begin{lemma}[Conditional $W_2$-Lipschitzness; the injection noise cancels]
\label{lem:condlip}
Under Assumptions~\ref{ass:compact} and \ref{ass:lip}, for all $\bth,\bth'\in\Om_\Theta$ and all
$\bm u,\bm u'\in\Om_H$,
\[
  W_2\big(Q_H(\cdot\mid\bth),Q_H(\cdot\mid\bth')\big)\le L_H\norm{\bth-\bth'},\qquad
  W_2\big(Q_P(\cdot\mid\bm u),Q_P(\cdot\mid\bm u')\big)\le L_P\norm{\bm u-\bm u'}.
\]
\end{lemma}
\begin{proof}
Couple $Q_H(\cdot\mid\bth)$ and $Q_H(\cdot\mid\bth')$ by drawing the \emph{same} latent $\bm z_H$ and
the \emph{same} injection noise $\bm\xi_H$: set $\bm U=\mathcal B_H(g_\phi(\bth,\bm z_H)+\sigma_H\bm\xi_H)$
and $\bm U'=\mathcal B_H(g_\phi(\bth',\bm z_H)+\sigma_H\bm\xi_H)$, which have the correct marginals, so
$(\bm U,\bm U')$ is a coupling. Since $\mathcal B_H$ is $1$-Lipschitz (Fact~\ref{fact:proj}) and the
shared term $\sigma_H\bm\xi_H$ cancels inside the difference,
\[
  \norm{\bm U-\bm U'}\le\big\|(g_\phi(\bth,\bm z_H)+\sigma_H\bm\xi_H)-(g_\phi(\bth',\bm z_H)+\sigma_H\bm\xi_H)\big\|
  =\norm{g_\phi(\bth,\bm z_H)-g_\phi(\bth',\bm z_H)}\le L_H\norm{\bth-\bth'}
\]
pathwise (Assumption~\ref{ass:lip}). By Fact~\ref{fact:w2basic},
$W_2\le(\E\norm{\bm U-\bm U'}^2)^{1/2}\le L_H\norm{\bth-\bth'}$. The $Q_P$ bound is identical. The
cancellation of $\sigma_H\bm\xi_H$ is the formal content of ``N4 sets the \emph{spread} of $Q_H$ but not
its \emph{sensitivity to the condition}.''
\end{proof}

\begin{lemma}[Sweep contraction]
\label{lem:sweep}
Let $K$ be the $\bth$-marginal one-sweep kernel,
$K(\bth,\cdot)=\int_{\Om_H}Q_P(\cdot\mid\bm u)\,Q_H(d\bm u\mid\bth)$ (draw $\bm u\sim Q_H(\cdot\mid\bth)$,
then $\bth'\sim Q_P(\cdot\mid\bm u)$). Under Assumptions~\ref{ass:compact},\ref{ass:cont},\ref{ass:lip},
for all $\bth,\bth'\in\Om_\Theta$,
\[
  W_2\big(K(\bth,\cdot),K(\bth',\cdot)\big)\le\kappa\,\norm{\bth-\bth'},\qquad \kappa:=L_HL_P .
\]
If moreover $\kappa<1$, then $K$ has a unique invariant $\nu^\ast\in\cP(\Om_\Theta)$ with
$W_2(\mu_0K^n,\nu^\ast)\le\kappa^n W_2(\mu_0,\nu^\ast)$, and $\nu^\ast$ is the $\bth$-marginal of the
joint invariant law $\pi^\ast$ of Theorem~\ref{thm:A}.
\end{lemma}
\begin{proof}
% [Rev.3-1] Proof reorganized into labelled steps (cosmetic; the
% argument is unchanged). [Rev.3-3] In Step 2 the operation is now
% called MIXING (against an optimal coupling), not gluing.
\emph{Step 1 (pathwise contraction).} Couple $K(\bth,\cdot),K(\bth',\cdot)$ by sharing all four
noises $(\bm z_H,\bm\xi_H,\bm z_P,\bm\xi_P)$. With $\bm U,\bm U'$ as in Lemma~\ref{lem:condlip} and
$\bm\Theta'=\mathcal B_\Theta(g_\psi(\bm U,\bm z_P)+\sigma_P\bm\xi_P)$,
$\tilde{\bm\Theta}'=\mathcal B_\Theta(g_\psi(\bm U',\bm z_P)+\sigma_P\bm\xi_P)$, the same cancellation
($\mathcal B_\Theta$ $1$-Lipschitz, shared $\sigma_P\bm\xi_P$) gives, pathwise,
\[
  \norm{\bm\Theta'-\tilde{\bm\Theta}'}\le\norm{g_\psi(\bm U,\bm z_P)-g_\psi(\bm U',\bm z_P)}
  \le L_P\norm{\bm U-\bm U'}\le L_PL_H\norm{\bth-\bth'}.
\]
Hence $W_2(K(\bth,\cdot),K(\bth',\cdot))\le(\E\norm{\bm\Theta'-\tilde{\bm\Theta}'}^2)^{1/2}\le\kappa\norm{\bth-\bth'}$.
($K$ is a Markov kernel on $\Om_\Theta$ by Fact~\ref{fact:compose}, exactly as for $T$, so
Fact~\ref{fact:banachw2} applies below.)

\emph{Step 2 (measure action).} The shared-noise construction of Step 1 defines an \emph{explicit}
coupling kernel $(\bth,\bth')\mapsto\gamma_{\bth,\bth'}:=\mathrm{law}(\bm\Theta',\tilde{\bm\Theta}')$,
jointly measurable in $(\bth,\bth')$ (a pushforward of the fixed noise law under a map that is jointly
continuous by Assumption~\ref{ass:cont} and Fact~\ref{fact:proj}); no measurable selection is needed
here. Mixing this kernel against an optimal coupling $\gamma\in\Gamma(\mu,\mu')$ of the two fixed
measures (exactly as in Fact~\ref{fact:msel}(iii), with the explicit kernel in place of a selected
one) and integrating the pathwise bound of Step 1 gives
$W_2(\mu K,\mu'K)^2\le\int\kappa^2\norm{\bth-\bth'}^2\,\gamma(d\bth,d\bth')=\kappa^2\,W_2(\mu,\mu')^2$.

\emph{Step 3 (invariant law and its identification).} When $\kappa<1$, Fact~\ref{fact:banachw2}
yields $\nu^\ast$ and the rate. Finally, if $(\bm u,\bth)\sim\pi^\ast$ then the swept
$\bth'\sim K(\bth,\cdot)$ has the same law as the $\bth$-marginal of $\pi^\ast$ (invariance of
$\pi^\ast$ under $T$), so that marginal is $K$-invariant and hence equals the unique $\nu^\ast$.
\end{proof}

\begin{remark}[Roles of the four noises; the collapse ``paradox'' dissolves]
\label{rem:noises}
The contraction constant $\kappa=L_HL_P$ is governed \emph{solely} by the condition-Lipschitz
constants (N1 via \cite{Bishop}); the injection noise $\sigma_H,\sigma_P$ (N4) and the latent (N3)
\emph{cancel} in the couplings above and do not enter $\kappa$. Conversely N4 is exactly what keeps
each $Q$ non-degenerate (full support, Lemma~\ref{lem:minor}) and sets the width of $\pi^\ast$; a
deterministic map ($\sigma\to0$) would $W_2$-contract to a \emph{point mass} --- the very mode
collapse we must avoid. Thus tightening the contraction (raise N1, lower $\kappa$) and preventing
collapse (keep N4$>0$) act on \emph{disjoint} parts of the analysis: there is no conflict, and the
apparent ``more contraction $\Rightarrow$ more collapse'' tension was an artifact of conflating N1
with N4.
\end{remark}

\begin{remark}[Honest logical dependency; existence vs.\ rate]
\label{rem:depend}
Under $\kappa<1$, Fact~\ref{fact:banachw2} already yields existence/uniqueness of $\nu^\ast$; so the
contraction, \emph{if it holds}, subsumes existence. The distinct value of Theorem~\ref{thm:A} is that
it proves existence/uniqueness of $\pi^\ast$ \emph{without} assuming any contraction --- from
$\sigma>0$ and compactness alone, quantities the algorithm guarantees. In the split we adopt,
Theorem~\ref{thm:A} (TV, minorization) supplies existence unconditionally, while
Lemma~\ref{lem:sweep} ($W_2$, condition-Lipschitz) supplies the \emph{rate} and, next, the stability
constant $1/(1-\kappa)$ for the steering bound (Theorem~B). We work on the $\bth$-marginal chain
$K$ because $p(\bth\mid\xobs)$ is the reported target; a joint contraction holds as well under a
suitably weighted product metric, with constant still governed by $L_H,L_P$, and is omitted.
\end{remark}

% ====================================================================
\section{Fixed-point perturbation}
\label{sec:perturb}

The rate constant $\kappa$ of Lemma~\ref{lem:sweep} now buys a \emph{stability} statement: two
contractive sweeps that are close as operators have close invariant laws.

\begin{lemma}[Fixed-point perturbation in $W_2$]
\label{lem:perturb}
Let $K$ be a Markov kernel on $\Om_\Theta$ whose measure action is a $\kappa$-contraction
($\kappa\in[0,1)$) with invariant $\nu^\ast$, and let $K'$ be any Markov kernel with an invariant
$\nu'$. Then
\[
  W_2(\nu^\ast,\nu')\ \le\ \frac{1}{1-\kappa}\,\sup_{\bth\in\Om_\Theta}\,W_2\big(K(\bth,\cdot),K'(\bth,\cdot)\big).
\]
\end{lemma}
\begin{proof}
Using $\nu^\ast=\nu^\ast K$ and $\nu'=\nu'K'$,
\[
  W_2(\nu^\ast,\nu')=W_2(\nu^\ast K,\nu'K')\le W_2(\nu^\ast K,\nu'K)+W_2(\nu'K,\nu'K').
\]
The first term is $\le\kappa\,W_2(\nu^\ast,\nu')$ ($\kappa$-contraction of $K$). For the second,
% [Rev.3-3] Terminology: the family of optimal couplings is MIXED against
% nu' (Fact msel(iii)); "glued" removed. Substance unchanged from Rev.2.
$W_2(\nu'K,\nu'K')\le\sup_{\bth} W_2(K(\bth,\cdot),K'(\bth,\cdot))$ by the averaged-kernel bound of
Fact~\ref{fact:msel}(iii) (a measurable family of optimal couplings, selected by
Fact~\ref{fact:msel}(ii) and mixed against $\nu'$). Rearranging gives the claim (all quantities are
finite since $\Om_\Theta$ is compact). Only $K$ need contract; $K'$ need only possess an invariant
law.
\end{proof}

% ====================================================================
\section{Theorem B: steering / consistency of the invariant law}
\label{sec:thmB}

% [Rev.3-5] Assumption suptrue is kept HERE (not moved into Sec. 6) so
% that Secs. 1--9 remain verbatim; it is standing for this section, and
% is now cited explicitly by Lemma opdisc and Theorem B (it was silently
% used in Rev.2: domain of the true sweep, definition of eps_H/eps_P,
% and coverage of the Lipschitz domains of Assumption lip all need it).
\begin{assumption}[The domain contains the true support]
\label{ass:suptrue}
$\operatorname{supp}\Pi^\dagger\subseteq\Om=\Om_H\times\Om_\Theta$, where $\Pi^\dagger:=p(\xhid,\bth\mid\xobs)$
is the true joint posterior. (Physically, the projection box is chosen to contain the plausible
parameter/state range.) Fix, once and for all, regular versions of the true conditionals
$p^\dagger_H(\cdot\mid\bth),p^\dagger_P(\cdot\mid\bm u)$ defined for every $\bth\in\Om_\Theta$,
$\bm u\in\Om_H$; being regular versions, these are Markov kernels.
\end{assumption}

% ---------------------------------------------------------------------
% [Rev.3-4] Target paragraph rewritten. WHY: (a) Rev.2 said "Pi-dagger is
% K-dagger-invariant", a category error (K-dagger acts on P(Om_Theta),
% not on joint laws on Om_H x Om_Theta); (b) the needed statement
% nu-dagger K-dagger = nu-dagger was carried by the slogan "Gibbs
% sampling is exact" without argument. HOW: the kernel-ness of K-dagger
% is pointed to Fact compose (as for T), and the invariance is PROVED by
% the same sweep-reproduction argument used in Prop. inc.
% ---------------------------------------------------------------------
\paragraph{The target.}
With the fixed versions $p^\dagger_H,p^\dagger_P$ of Assumption~\ref{ass:suptrue}, define the true
sweep
\[
  K^\dagger(\bth,\cdot)\ :=\ \int p^\dagger_P(\cdot\mid\bm u)\,p^\dagger_H(d\bm u\mid\bth).
\]
This is a Markov kernel on $\Om_\Theta$: Assumption~\ref{ass:suptrue} keeps every step inside $\Om$,
and measurability follows from Fact~\ref{fact:compose} exactly as for $T$ (Lemma~\ref{lem:measurable}).
Moreover the forward sweep \emph{reproduces} $\Pi^\dagger$: if $(\bm u,\bth)\sim\Pi^\dagger$ and one
redraws $\bm u'\sim p^\dagger_H(\cdot\mid\bth)$, then $(\bm u',\bth)\sim\Pi^\dagger$ (as $p^\dagger_H$
is its $\bm u\mid\bth$ conditional); redrawing $\bth'\sim p^\dagger_P(\cdot\mid\bm u')$ then gives
$(\bm u',\bth')\sim\Pi^\dagger$ (as $p^\dagger_P$ is its $\bth\mid\bm u$ conditional). Taking
$\bth$-marginals of the first and last joints yields
\[
  \nu^\dagger K^\dagger=\nu^\dagger,\qquad \nu^\dagger:=p(\bth\mid\xobs)\ \text{the $\bth$-marginal of }\Pi^\dagger .
\]
This $\nu^\dagger$ is the estimand.

\paragraph{Approximation defects.}
Define
\[
  \varepsilon_H:=\sup_{\bth\in\Om_\Theta} W_2\big(Q_H(\cdot\mid\bth),\,p^\dagger_H(\cdot\mid\bth)\big),\qquad
  \varepsilon_P:=\sup_{\bm u\in\Om_H} W_2\big(Q_P(\cdot\mid\bm u),\,p^\dagger_P(\cdot\mid\bm u)\big),
\]
the worst-case $W_2$ gaps between each \emph{learned} conditional and the corresponding \emph{true}
conditional, both suprema taken over the compact domains and with respect to the fixed versions of
Assumption~\ref{ass:suptrue}. (These absorb finite-sample/finite-capacity error, the injection-noise
smoothing, and the boundary projection; see Remark~\ref{rem:edist}.)

% [Rev.3-5] Assumption citation now includes suptrue.
\begin{lemma}[Operator discrepancy]
\label{lem:opdisc}
Under Assumptions~\ref{ass:compact}--\ref{ass:lip} and~\ref{ass:suptrue},
$\ \sup_{\bth\in\Om_\Theta} W_2\big(K(\bth,\cdot),K^\dagger(\bth,\cdot)\big)\le L_P\,\varepsilon_H+\varepsilon_P.$
\end{lemma}
\begin{proof}
Interpolate through $K'(\bth,\cdot)=\int Q_P(\cdot\mid\bm u)\,p^\dagger_H(d\bm u\mid\bth)$ (true
$H$-step, learned $P$-step):
\[
  W_2\big(K(\bth,\cdot),K^\dagger(\bth,\cdot)\big)\le W_2\big(K(\bth,\cdot),K'(\bth,\cdot)\big)+W_2\big(K'(\bth,\cdot),K^\dagger(\bth,\cdot)\big).
\]

\emph{First term (swap $H$).} At the fixed $\bth$, take one optimal coupling $(\bm U,\bm U^\dagger)$
of $Q_H(\cdot\mid\bth),p^\dagger_H(\cdot\mid\bth)$ with
$\E\norm{\bm U-\bm U^\dagger}^2=W_2\big(Q_H(\cdot\mid\bth),p^\dagger_H(\cdot\mid\bth)\big)^2$; note
$\bm U,\bm U^\dagger\in\Om_H$ a.s.\ ($Q_H$ is post-projection; $p^\dagger_H$ by
Assumption~\ref{ass:suptrue}), so the Lipschitz domain of Assumption~\ref{ass:lip} covers both. Then
draw $\bth',\bth'^\dagger$ from $Q_P(\cdot\mid\bm U),Q_P(\cdot\mid\bm U^\dagger)$ by shared noise
$(\bm z_P,\bm\xi_P)$ \emph{independent of} $(\bm U,\bm U^\dagger)$ (Lemma~\ref{lem:condlip}), so
$\norm{\bth'-\bth'^\dagger}\le L_P\norm{\bm U-\bm U^\dagger}$. As this is a single fixed coupling, no
selection is needed;
$W_2(K,K')\le L_P\,W_2\big(Q_H(\cdot\mid\bth),p^\dagger_H(\cdot\mid\bth)\big)\le L_P\varepsilon_H$.

\emph{Second term (swap $P$).} Both use $p^\dagger_H(\cdot\mid\bth)$ for the $\bm u$-step; couple with
the \emph{same} $\bm u$, then couple $Q_P(\cdot\mid\bm u),p^\dagger_P(\cdot\mid\bm u)$ by a
\emph{measurable} family of optimal couplings $\bm u\mapsto\gamma_{\bm u}$
(Fact~\ref{fact:msel}(ii); both are Markov kernels --- $Q_P$ by Lemma~\ref{lem:measurable},
$p^\dagger_P$ by the fixed regular version of Assumption~\ref{ass:suptrue}) and mix against
$p^\dagger_H(\cdot\mid\bth)$ (Fact~\ref{fact:msel}(iii)), giving
$W_2(K',K^\dagger)\le\big(\int W_2\big(Q_P(\cdot\mid\bm u),p^\dagger_P(\cdot\mid\bm u)\big)^2\,p^\dagger_H(d\bm u\mid\bth)\big)^{1/2}\le\varepsilon_P$.

Sum and take $\sup_{\bth}$.
\end{proof}

% [Rev.3-5] Assumption citation now includes suptrue.
\begin{theorem}[Steering bound]
\label{thm:B}
Under Assumptions~\ref{ass:compact}--\ref{ass:lip} and~\ref{ass:suptrue}, with $\kappa=L_HL_P<1$, the
invariant $\bth$-law $\nu^\ast$ of the learned sweep satisfies
\[
  W_2\big(\nu^\ast,\ p(\bth\mid\xobs)\big)\ \le\ \frac{L_P\,\varepsilon_H+\varepsilon_P}{1-\kappa}.
\]
In particular, if $\varepsilon_H=\varepsilon_P=0$ then $\nu^\ast=p(\bth\mid\xobs)$ exactly
(\textbf{Corollary B1}, the well-specified/compatible case).
\end{theorem}
\begin{proof}
$K$ is a $\kappa$-contraction with invariant $\nu^\ast$ (Lemma~\ref{lem:sweep}) and $K^\dagger$ is a
Markov kernel on $\Om_\Theta$ with invariant $\nu^\dagger=p(\bth\mid\xobs)$ (the target paragraph
above). Apply Lemma~\ref{lem:perturb} and bound the operator discrepancy by Lemma~\ref{lem:opdisc}.
\end{proof}

\paragraph{The computable self-consistency defect.}
Theorem~\ref{thm:B} bounds the distance to the \emph{true} posterior, but
$\varepsilon_H,\varepsilon_P$ require the unknown truth. We therefore also define an intrinsic,
\emph{model-only} defect that certifies whether the learned pair even admits a coherent joint. Let
$\mu^\ast(d\bm u):=\int Q_H(d\bm u\mid\bth)\,\nu^\ast(d\bth)$ be the $\bm u$-marginal the forward
sweep induces, and form the two joints on $\Om_\Theta\times\Om_H$
\[
  \Pi^{\rightarrow}(d\bth,d\bm u):=\nu^\ast(d\bth)\,Q_H(d\bm u\mid\bth),\qquad
  \Pi^{\leftarrow}(d\bth,d\bm u):=Q_P(d\bth\mid\bm u)\,\mu^\ast(d\bm u),\qquad
  \varepsilon_{\mathrm{inc}}:=W_2\big(\Pi^{\rightarrow},\Pi^{\leftarrow}\big).
\]
Both have the same $\bm u$-marginal $\mu^\ast$; they differ only in the $\bth\mid\bm u$ law (the
forward posterior of $\bth$ under $\Pi^\rightarrow$ vs.\ the learned $Q_P$), so
$\varepsilon_{\mathrm{inc}}$ is the ``forward-vs-backward sweep inconsistency'' and is estimable from
samples of the trained model alone.

% [Rev.3-6] Standing assumptions of the Proposition made explicit
% (uniqueness of nu* is via Lemma sweep, which needs Assumptions 1--3).
\begin{proposition}[$\varepsilon_{\mathrm{inc}}$ certifies compatibility]
\label{prop:inc}
Let Assumptions~\ref{ass:compact}--\ref{ass:lip} hold with $\kappa=L_HL_P<1$. Then
$\varepsilon_{\mathrm{inc}}=0$ iff $Q_H,Q_P$ are compatible (equal, a.e., the conditionals of a
common joint); in that case the witnessing joint is $\Pi=\Pi^\rightarrow=\Pi^\leftarrow$, and
$\nu^\ast$ is exactly its $\bth$-marginal --- i.e.\ the sampler draws from a genuine joint's marginal.
\end{proposition}
\begin{proof}
$(\Leftarrow)$ Suppose a joint $\Pi'$ has $\bm u\mid\bth$-conditional $Q_H$ and
$\bth\mid\bm u$-conditional $Q_P$ (a.e.). Running one forward sweep from $\Pi'$ reproduces $\Pi'$, so
its $\bth$-marginal $\Pi'_{\bth}$ is $K$-invariant, hence $\Pi'_{\bth}=\nu^\ast$ (unique, $\kappa<1$).
Thus $\Pi'=\nu^\ast\otimes Q_H=\Pi^\rightarrow$; its $\bm u$-marginal is $\mu^\ast$ and its
$\bth\mid\bm u$-conditional is $Q_P$, so also $\Pi'=Q_P\otimes\mu^\ast=\Pi^\leftarrow$, giving
$\varepsilon_{\mathrm{inc}}=0$.

$(\Rightarrow)$ If $\varepsilon_{\mathrm{inc}}=0$ then $\Pi^\rightarrow=\Pi^\leftarrow=:\Pi$; by
construction $\Pi$ has $\bm u\mid\bth$-conditional $Q_H$ (from $\Pi^\rightarrow$) and
$\bth\mid\bm u$-conditional $Q_P$ (from $\Pi^\leftarrow$), so $Q_H,Q_P$ are compatible with witness
$\Pi$, whose $\bth$-marginal is $\nu^\ast$.
\end{proof}

% [Rev.3-6] kappa<1 noted as in force (nu*, mu* are otherwise undefined).
\begin{remark}[$\Pi^\leftarrow$ is the stationary joint]
\label{rem:pileftpi}
(Throughout this remark $\kappa<1$ is in force, so $\nu^\ast$ --- and with it $\mu^\ast$ --- is
defined.) The backward-assembled joint is exactly the joint invariant law: at stationarity the swept
state $(\bm u',\bth')$ has $\bm u'\sim\mu^\ast$ and $\bth'\sim Q_P(\cdot\mid\bm u')$, i.e.\
$\pi^\ast=\Pi^\leftarrow$ (up to the coordinate identification
$\Om_H\times\Om_\Theta\cong\Om_\Theta\times\Om_H$). Hence
$\varepsilon_{\mathrm{inc}}=W_2(\Pi^\rightarrow,\pi^\ast)$: the self-consistency defect measures how
far the forward pairing $(\bth,\bm u\!\sim\!Q_H(\cdot\mid\bth))$ sits from the sampler's own
stationary joint --- a quantity one estimates directly from a single run.
\end{remark}

% [Rev.3-8] Hedge sharpened: the swap-learned-first triangle route
% suggests the constant closes with kappa, L_H, L_P only (no regularity
% of the TRUE conditionals is needed, since Lemma condlip is only ever
% applied to the LEARNED conditionals); precise constant still deferred.
\begin{remark}[Why the defects are not summed --- no double counting]
\label{rem:nodouble}
It is tempting to add $\varepsilon_{\mathrm{inc}}$ to the bound of Theorem~\ref{thm:B}; this would
double count. Since the true conditionals are compatible, the true joints built forward and backward
coincide ($\Pi^\rightarrow_\dagger=\Pi^\leftarrow_\dagger=\Pi^\dagger$), so by a triangle inequality
$\varepsilon_{\mathrm{inc}}\le W_2(\Pi^\rightarrow,\Pi^\dagger)+W_2(\Pi^\dagger,\Pi^\leftarrow)$. Each
term can in turn be controlled by swapping the \emph{learned} conditional first --- e.g.\
$W_2\big(Q_H(\cdot\mid\bth),p^\dagger_H(\cdot\mid\bth^\dagger)\big)\le\varepsilon_H+L_H\norm{\bth-\bth^\dagger}$
by the triangle inequality and Lemma~\ref{lem:condlip} applied to the learned $Q_H$ --- together with
$W_2(\nu^\ast,\nu^\dagger)$ from Theorem~\ref{thm:B}. One expects, accordingly, a bound
$\varepsilon_{\mathrm{inc}}\le C(\kappa,L_H,L_P)\,(\varepsilon_H+\varepsilon_P)$; since
Lemma~\ref{lem:condlip} is only ever invoked on the learned conditionals along this route, no
regularity of the \emph{true} conditionals should be needed. We defer the precise constant to a later
note. The qualitative point stands regardless: $\varepsilon_{\mathrm{inc}}$ is \emph{dominated by} the
approximation defects, not independent of them --- $\varepsilon_H,\varepsilon_P$ govern the distance
to truth (Theorem~\ref{thm:B}), while $\varepsilon_{\mathrm{inc}}$ is its computable \emph{shadow},
the quantity one can actually monitor (e.g.\ across noise scales; cf.\ the ablation program).
\end{remark}

% [Rev.3-7] Law of Y specified (stationary PRE-projection proposal), and
% sqrt(d) dimension factors flagged (consistent with Remark rem:scope's
% policy of flagging dimension dependence).
\begin{remark}[Where the distortion defect lives; the $\sigma$ decomposition]
\label{rem:edist}
Here $p^\dagger_H,p^\dagger_P$ are the \emph{clean} true conditionals, so $\varepsilon_H,\varepsilon_P$
silently include the injection-noise smoothing (N4) and the boundary projection --- i.e.\ a distortion
defect
$\varepsilon_{\mathrm{dist}}=O\big((\sigma_H+\sigma_P)\sqrt{d}\big)+\big(\E\norm{\bm Y-\mathcal B(\bm Y)}^2\big)^{1/2}$
is folded in, where $\bm Y$ denotes the stationary \emph{pre-projection} proposal (for the $H$-step,
$\bm Y=g_\phi(\bth,\bm z_H)+\sigma_H\bm\xi_H$ with $\bth\sim\nu^\ast$; analogously for the $P$-step,
with the two contributions combined) and $\mathcal B$ the corresponding projection. The last term is
an RMS projection \emph{displacement} --- a length, not the boundary \emph{mass}
$\pi^\ast(\partial\Om)$; the two are distinct. The $\sqrt d$ factors make the dimension dependence
explicit (adding $\Normal(\bm0,\sigma^2I_d)$ noise moves a law by $\sigma\sqrt d$ in $W_2$; cf.\ the
dimension-flagging policy of Remark~\ref{rem:scope}). To display the distortion separately, reference
instead the $\sigma$-smoothed, projected true conditionals $p^{\dagger,\sigma}_H,p^{\dagger,\sigma}_P$:
then $\varepsilon_H,\varepsilon_P$ reduce to pure learning error (vanishing with data/capacity), and
one adds the single term $W_2(\nu^{\dagger,\sigma},\nu^\dagger)=O(\sigma\sqrt d)$ for the smoothing
bias. Because Theorem~\ref{thm:A} and non-collapse require $\sigma_H,\sigma_P>0$ (N4), this term has
an irreducible floor; in the non-identifiable regime it is small relative to the width of the true
posterior ridge, so the ridge is faithfully covered.
\end{remark}

% ====================================================================
% [Rev.3-11] Bibliography parentheticals extended: Villani Cor. 5.22
% (measurable selection) and Kallenberg (measurable structure of P(Om)),
% both now load-bearing via Facts kernmeas/msel.
\begin{thebibliography}{9}
\bibitem{Bishop} C.~M.~Bishop. Training with noise is equivalent to Tikhonov regularization.
  \emph{Neural Computation}, 7(1):108--116, 1995.
\bibitem{Folland} G.~B.~Folland. \emph{Real Analysis: Modern Techniques and Their Applications},
  2nd ed. Wiley, 1999. (Radon--Nikodym, Ch.~3; $L^p$ completeness, Ch.~6.)
\bibitem{HiriartLemarechal} J.-B.~Hiriart-Urruty and C.~Lemar\'echal.
  \emph{Fundamentals of Convex Analysis}. Springer, 2001. (Projection onto a convex set, Ch.~III.)
\bibitem{Kallenberg} O.~Kallenberg. \emph{Foundations of Modern Probability}, 2nd ed. Springer, 2002.
  (Kernels and section integration, Ch.~1; measurable structure of $\cP(\Om)$ --- evaluation
  $\sigma$-algebra vs.\ weak Borel --- Chs.~1 and 4; Ionescu--Tulcea, Ch.~6. Chapter numbers to be
  verified against the edition in use.)
\bibitem{MeynTweedie} S.~Meyn and R.~L.~Tweedie.
  \emph{Markov Chains and Stochastic Stability}, 2nd ed. Cambridge Univ.\ Press, 2009.
  (Doeblin/uniform ergodicity, Ch.~16.)
\bibitem{Villani} C.~Villani. \emph{Optimal Transport: Old and New}. Springer, 2009.
  ($W_2$ metric and its properties, Ch.~6; gluing lemma, Ch.~1; measurable selection of optimal
  plans, Cor.~5.22.)
\end{thebibliography}

\end{document}